\documentclass[10pt,a4paper]{amsart}
\usepackage[margin=19mm]{geometry}
\usepackage{amsmath,amssymb,mathtools}
\usepackage[scr=rsfs,cal=euler]{mathalfa}
\usepackage{xfrac}
\usepackage[unicode,hidelinks,bookmarks=false]{hyperref}
\newcommand{\ie}{i.e.\ }
\newcommand{\cf}{cf.\ }
\newcommand{\resp}{respectively\ }
\newcommand{\Subsection}{\subsection}
\providecommand{\nobreakdash}{\nobreak\hbox{-}}
\setlength{\emergencystretch}{2em}
\multlinegap0pt
\usepackage{tikz,aliascnt}
\usetikzlibrary{cd,calc,arrows}
\usepackage[capitalize]{cleveref}
\newcommand{\skpt}{\par\smallskip}
\newtheorem{theorem}{Theorem}[section]
\newaliascnt{lemma}{theorem}
\newtheorem{lemma}[lemma]{Lemma}
\aliascntresetthe{lemma}
\crefname{lemma}{Lemma}{Lemmas}
\newaliascnt{proposition}{theorem}
\newtheorem{proposition}[proposition]{Proposition}
\aliascntresetthe{proposition}
\crefname{proposition}{Proposition}{Propositions}
\newaliascnt{claim}{theorem}
\newtheorem{claim}[claim]{Claim}
\aliascntresetthe{claim}
\crefname{claim}{Claim}{Claims}
\newaliascnt{corollary}{theorem}
\newtheorem{corollary}[corollary]{Corollary}
\aliascntresetthe{corollary}
\crefname{corollary}{Corollary}{Corollarys}
\theoremstyle{definition}
\newaliascnt{definition}{theorem}
\newtheorem{definition}[definition]{Definition}
\aliascntresetthe{definition}
\crefname{definition}{Definition}{Definitions}
\newaliascnt{example}{theorem}
\newtheorem{example}[example]{Example}
\aliascntresetthe{example}
\crefname{example}{Example}{Examples}
\newaliascnt{remark}{theorem}
\newtheorem{remark}[remark]{Remark}
\aliascntresetthe{remark}
\crefname{remark}{Remark}{Remarks}
\newaliascnt{question}{theorem}
\newtheorem{question}[question]{Question}
\aliascntresetthe{question}
\crefname{question}{Question}{Questions}
\newtheorem{introdefinition}{Definition}
\renewcommand{\theintrodefinition}{\Alph{introdefinition}}
\crefname{introdefinition}{Definition}{Definitions}
\crefname{enumi}{}{}
\creflabelformat{enumi}{#2\textup{(#1)}#3}
\let\ts\textstyle
\newcommand\mto{\mathchoice{\longmapsto}{\mapsto}{\mapsto}{\mapsto}}
\newcommand\hto{\mathrel{\lhook\joinrel\to}}
\newcommand\from{\mathchoice{\longleftarrow}{\leftarrow}{\leftarrow}{\leftarrow}}
\newcommand\hfrom{\mathrel{\from\joinrel\rhook}}

\newcommand{\spfrac}[2]{\sfrac{#1}{(#2)}}

\newcommand{\Aut}{\mathrm{Aut}}
\newcommand{\LO}{\mathrm{LO}}
\newcommand{\Sym}{\mathrm{Sym}}
\newcommand{\Cyc}{\mathrm{Cyc}}
\newcommand{\Alt}{\mathrm{Alt}}
\newcommand{\Stab}{\mathrm{Stab}}
\newcommand{\Homeo}{\mathrm{Homeo}}
\newcommand{\Ends}{\mathrm{Ends}}
\newcommand{\Split}{\mathrm{Split}}
\newcommand{\Prob}{\mathrm{Prob}}
\newcommand{\A}{\mathcal{A}}
\newcommand{\B}{\mathcal{B}}
\newcommand{\R}{\mathcal{R}}
\newcommand{\U}{\mathcal{U}}
\newcommand{\K}{\mathcal{K}}
\newcommand{\Chi}{\mathcal{X}}
\newcommand{\Br}{\mathrm{Br}}
\newcommand{\interior}[1]{\overset{\circ}{#1}}
\newcommand{\C}{\mathbf{C}}

\title[Regular-rabbit universal homeomorphism group]{Homeomorphism groups of basilica, rabbit and airplane Julia sets: original Theorem1.29}
\author{Bruno Duchesne and Matteo Tarocchi}
\date{Source: J.\'E.P.--M.13(2026),255--320; DOI10.5802/jep.327}
\begin{document}\maketitle
\noindent\textit{Editorial scope.} The counted stable source slot is original Theorem1.29 for finite $n\ge3$. The full original general statement remains verbatim; the basilica and infinite-order cases receive no count. Complete circle-tree, topology, legal-coloring, lifting and continuity arguments are retained as uncounted core, together with the required cyclic-order appendix. Illustrative Julia images and finite coloring pictures, and later group applications, are omitted. No author mathematical expression or assumption is repaired.
Recall that a \emph{Peano continuum} is a compact, metrizable, connected, and locally connected space.
Peano continua are path-connected, for example because of \cite[Th.\,8.23]{continua}. By~an \emph{arc} of such a continuum, we~mean a subset that is homeomorphic to the topological interval $[0,1]$ and by a \emph{circle} we mean a subset that is homeomorphic to the topological circle $\mathbf{S}^1$.
Finally, recall that the \emph{order} of a point is the number of connected components of its complement, and a point is a \emph{cut point} if its order is at least $2$.
With this in mind, we~present our abstract definition of rabbits.

\begin{introdefinition}
\label{def.rabbit}
A \emph{rabbit} is a Peano continuum that enjoys the following properties.
\begin{enumerate}
\item\label{def.rabbit1} The set of points that belong to multiple circles is dense.
\item\label{def.rabbit2} If two points belong to distinct circles, they are separated by a cut point.
\item\label{def.rabbit3} For any two points each of which belongs to some circle, every arc between them is included in a union of finitely many circles.
\end{enumerate}
A rabbit is $n$-\emph{regular} ($n\geq2$) if its cut points all have the same order $n \in \mathbf{N}_{\geq 2} \cup \{\infty\}$. A~\emph{basilica} is a $2$-regular rabbit and a \emph{Douady rabbit} is a $3$-regular rabbit.
\end{introdefinition}

\section{The rabbits and the basilica}
\label{sec.rabbits}


\subsection{Topological characterization of rabbits}

Throughout this work, given points $a,b \in X$ (or subsets $A,B \subseteq X$), we~say that a point $p$ or a subset $S \subset X$ \emph{separates} $a$ and $b$ if they lie in distinct connected components of $X \setminus \{p\}$ or $X \setminus S$ (or~separates $A$ and $B$ if they are included in distinct connected components of $X \setminus \{p\}$ or $X \setminus S$).

Let us now focus on the family of rabbits, which we defined in \cref{def.rabbit}.


\setcounter{theorem}{1}
\begin{definition}
\label{def.rabbit.points}
A point of a rabbit is said to be:
\begin{itemize}
\item an \emph{end point} if it does not belong to any circle,
\item a \emph{cut point} if its complement is not connected,
\item a \emph{regular point} if it belongs to a circle and is not a cut point (equivalently thanks to \cref{lem.circles.intersection} right below, if it belongs to a unique circle).
\end{itemize}
\end{definition}

The reason behind the name of end points will be clarified with \cref{rmk.rabbit.end.points}, after the introduction of the tree of circles.

The next two lemmas show why all conditions of \cref{def.rabbit} are essential.

\begin{lemma}
\label{lem.circles.intersection}
The intersection between two distinct circles of a rabbit is either empty or consists of a sole cut point.
\end{lemma}

\begin{proof}
Consider two distinct circles $C_1$ and $C_2$ of a rabbit. It~is impossible that one includes the other, so~there exist points $p \in C_1 \setminus C_2$ and $q \in C_2 \setminus C_1$. By~condition~\cref{def.rabbit2} of \cref{def.rabbit}, $p$ and $q$ are separated by a cut point.
Because of this, $C_1$ and $C_2$ cannot meet at more than one point, as otherwise $p$ and $q$ could not be separated by a cut point.
Then $C_1 \cap C_2$ is either empty or consists of a sole cut point.
\end{proof}

\begin{lemma}
\label{lem.cut.points.order}
The order of a cut point is equal to the number of circles that include~it.
\end{lemma}

\begin{proof}
We need to show that the removal of a cut point $p$ from a rabbit $\R$ produces a distinct connected component for each circle $C$ that contains $p$.

By condition \cref{def.rabbit2} of \cref{def.rabbit} there is a distinct connected component for each circle that meets at $p$.
Indeed, if two circles $C_1$ and $C_2$ at $p$ are such that $C_1 \setminus \{p\}$ and $C_2 \setminus \{p\}$ lie in the same connected component of $\R \setminus \{p\}$, then it would be impossible to separate $C_1$ and $C_2$ with a single cut point in $\R$, which contradicts condition~\cref{def.rabbit2}.

For the converse, if $A$ is a connected component of $\R \setminus \{p\}$, we~need to show that $A \cup \{p\}$ meets $p$ in a circle.
Fix another connected component $B$ at $p$, which is going to be useful to apply condition \cref{def.rabbit3} of \cref{def.rabbit}.
Since $\R$ is locally connected (by~the definition of Peano continua), $A$ and $B$ are open, so~there exist points $a \in A$ and $b \in B$ that belong to multiple circles by condition \cref{def.rabbit1} of \cref{def.rabbit}.
Fix any arc $[a,b]$ and note that it must contain $p$, since $p$ separates $A \ni a$ and $B \ni b$. By~condition~\cref{def.rabbit3} of \cref{def.rabbit}, since both $a$ and $b$ belong to some circle, the arc $[a,b]$ is included in a union of finitely many circles. In~particular, its subset $[a,p)$ is included in a union of finitely many circles.
Then there is some $q \in [a,p)$ such that $[q,p)$ is included in some circle $C$.
Since $C$ is closed, we~have that $p \in C$.
Finally, observe that if a circle intersects $A$ non-trivially then it must be entirely included in $A \cup \{p\}$.
Thus, $C \subseteq A$ and we are done.
\end{proof}

\begin{corollary}
\label{cor.rabbit.cut.points}
A point of a rabbit is a cut point if and only if it belongs to multiple circles.
\end{corollary}

\begin{remark}
\label{rmk.rabbit.countable}
Note that the order of a cut point of a rabbit is at most countably infinite.
Indeed, since rabbits are Peano continua they are locally connected second countable spaces, so~the complement of a point is a union of disjoint open connected subsets and, by second countability, there are at most countably many such subsets.

In a similar fashion, the removal of a circle must produce at most countably many connected components, so~each circle contains at most countably many cut points.
Because of condition~\cref{def.rabbit1} of \cref{def.rabbit}, it~must contain at least countably many, so~each circle has precisely countably many cut points.
\end{remark}

\begin{lemma}
\label{lem.rabbit.arcs.and.circles}
Any point of an arc minus its end points is contained in some circle.
\end{lemma}

\begin{proof}
Let $I = [a,b]$ be an arc in a rabbit $\R$ and let $p \in (a,b) \coloneqq [a,b] \setminus \{a,b\}$. If~$p$ is a cut point, then it belongs to multiple circles by \cref{cor.rabbit.cut.points}, so~there is nothing to prove. If~$p$ is not a cut point, then $\R \setminus \{p\}$ is connected and thus path-connected because $\R$ is a Peano continuum. It~follows that there exists an arc $I^*$ that joins $a$ and $b$ while avoiding $p$.
Then $I \cup I^*$ must include some circle that contains $p$.
\end{proof}

\begin{lemma}
The set of cut points of a rabbit is arcwise dense, \ie it has non-trivial intersection with any arc.
\end{lemma}

\begin{proof}
Let us show that any arc $I=[a,b]$ contains cut points.
Fix a point $p \in (a,b)$ and assume that it is not a cut point.
Then, up to taking a subinterval, we~can assume that $(a,b)$ is entirely included in some circle $C$.
Let $U$ be a path-connected neighborhood of $p$ that does not contain $a$ and $b$ (which can be found because $\R$ is locally path-connected). By~condition~\cref{def.rabbit1} of \cref{def.rabbit}, $U$ contains some point $p'$ that belongs to some circle.
Let $I'$ be an arc in $U$ that joins $p$ and $p'$.
The first point~$q$ in which $I'$ meets $C$ (coming from $p'$) must also belong to $I$, because $U$ does not contain $a$ and $b$.
The point $q$ must be a cut point, so~we are done.
\end{proof}

\begin{proposition}
\label{prop.subbasis}
Let $X$ be a Hausdorff compact space and consider a point $x \in X$.
Assume that $\mathcal{O}_x$ is a collection of open subsets of $X$ containing $x$ and with the property that, for every $y \in X$, there exists $U_y \in \mathcal{O}_x$ such that $y \notin \overline{U_y}$.
Then $\mathcal{O}_x$ is a subbasis of neighborhoods at $x$.

Moreover, given such collections $\mathcal{O}_x$ for every $x \in X$, their union $\mathcal{O} = \bigcup_{x \in X} \mathcal{O}_x$ is a subbasis for the topology of $X$.
\end{proposition}

\begin{proof}
Let $V$ be an open subset of $X$ and $x \in V$. We~want to show that there exist finitely many elements of the collection $\mathcal{O}_x$ whose intersection contains $x$ and is included in $V$.
Observe that
\[
 \bigcap \left\{ \overline{U_y} \mid y \in X \setminus V \right\} \cap (X \setminus V) = \emptyset, 
\]
otherwise there would exist $y \in X \setminus V$ contained in $\overline{U_y}$, which is a contradiction.
Since $X$ is compact and the infinite intersection in the previous equation is empty, there exist $y_1, \dots, y_n \in X \setminus V$ such that
\[
 \overline{U_{y_1}} \cap \dots \cap \overline{U_{y_n}} \cap (X \setminus V) = \emptyset, 
\]
meaning that $\overline{U_{y_1}} \cap \dots \cap \overline{U_{y_n}} \subseteq V$, as needed.
\end{proof}

Using this proposition, it~is easy to show the following useful fact.

\begin{corollary}
\label{cor.rabbit.subbasis}
Given a rabbit $\R$, the collection of all connected components of~$\R$ minus finitely many of its circles is a subbasis for the topology of $\R$.
\end{corollary}

\subsection{The tree of circles of a rabbit}\label{sec.tree.circles.rabbit}

In this subsection we will see how the circles of a rabbit are arranged in a tree-like manner, which is going to be essential when studying their homeomorphism groups.

\begin{definition}
\label{def.traverse}
Given an arc $I$ in a rabbit $\R$ and a subset $S \subseteq \R$, we~say that $I$ \emph{traverses} $S$ if the intersection $I \cap S$ includes a non-degenerate arc (\ie an arc that is not a singleton).
\end{definition}

We will often be using the previous definition in the case in which $S$ is a circle. In~that case, because of condition \cref{def.rabbit2} of \cref{def.rabbit}, we~can equivalently say that an arc $I$ traverses a circle $C$ if $I \cap C$ contains at least two points.

\begin{definition}
\label{def.tree.of.circles}
Given a rabbit $\R$, its \emph{tree of circles} $T_\R$ is the (undirected) graph defined as follows.
\begin{itemize}
\item The set of vertices consists of the circles and the cut points of $\R$.
\item There is an edge between a circle $C$ and a cut point $p$ when $p \in C$.
\end{itemize}
\end{definition}

The fact that the graph $T_\R$ is really a tree is proved in \cref{lem.tree.of.circles.is.tree}.

\begin{remark}
\label{rmk.arcs.and.paths}
Because of \cref{lem.circles.intersection} and by condition \cref{def.rabbit3} of \cref{def.rabbit}, there is a natural surjective map from the set of arcs of a rabbit $\R$ joining cut points or regular points onto the set of paths of the graph $T_\R$.
An arc $I = [p,q]$ is mapped to the sequence of circles traversed by $I$ (\cref{def.traverse}), interposed by the cut points between them. If~$p$ is a cut point, then the first vertex of the path is $p$ itself;
if $p$ is a regular point that belongs to a unique circle $C$, then the first vertex is $C$.
The same goes for $q$.
The preimages of this map are only distinguished by the choice of one of the two orientations of each traversed circle and possibly by a choice of regular point of $C$ when the initial (terminal) vertex of the path is a circle $C$.
\end{remark}

\begin{lemma}
\label{lem.tree.of.circles.is.tree}
Given a rabbit $\R$, the graph $T_\R$ is a tree.
Moreover, if a vertex is a circle then its degree is countably infinite and if it is instead a cut point then its degree is the order of the cut point.
\end{lemma}

\begin{proof}
Since rabbits are path-connected (because they are Peano continua), by the previous \cref{rmk.arcs.and.paths} the graph $T_\R$ is connected.
Now, a cycle in $T_\R$ would need to contain at least two distinct cut points and two distinct circles as vertices.
This would correspond to two distinct circles being joined by two arcs that travel through distinct cut points.
Then the two circles would not be separated by a cut point, which would violate condition \cref{def.rabbit2} of \cref{def.rabbit}.
Hence, $T_\R$ is a tree.

By \cref{lem.cut.points.order}, each cut point $p$ belongs to one circle for each of the connected components of $\R \setminus \{p\}$.
Thus, the degree of $p$ in $T_\R$ is the order of the cut point. As~noted in \cref{rmk.rabbit.countable}, each circle contains countably many cut points, so~its degree in~$T_\R$ is infinite.
\end{proof}

\begin{remark}
\label{rmk.rabbit.end.points}
The set of end points of rabbits (\cref{def.rabbit.points}) corresponds to the set $\partial T$ of ends of the tree of circles (\ie the set of semi-infinite paths from an arbitrary fixed vertex).

Indeed, if $e$ is an end point and $C_0$ is an arbitrary circle, consider an arc $I$ from $C_0$ to $e$. By~\cref{lem.rabbit.arcs.and.circles}, $I \setminus \{e\}$ is included in a union of circles. The arc $I$ thus traverses (\cref{def.traverse}) an infinite sequence of distinct circles $(C_1, C_2, \dots)$ which corresponds, by \cref{rmk.arcs.and.paths}, to a semi-infinite path in $T_\R$ from the vertex $C_0$.
Conversely, a semi-infinite path in $T_\R$ from $C_0$ corresponds to an infinite sequence of distinct circles $(C_1, C_2, \dots)$. If~$B_k$ is the unique connected component of $\R \setminus C_k$ that meets $C_{k+1}$, then the intersection of all the $B_k$'s is non-empty (because each finite intersection $B_0 \cap \dots \cap B_k$ is non-empty and $\R$ is compact).
Points in the intersection of all the $B_k$'s are end points, as they cannot belong to a circle by condition \cref{def.rabbit3} of \cref{def.rabbit}, so~the correspondence between end points and semi-infinite paths from $C_0$ is surjective.
Finally, given two distinct end points $e_1$ and $e_2$, since $\R$ is a metrizable space the arcs $I_1$ and $I_2$ from $C_0$ to $e_1$ and $e_2$, respectively, must eventually diverge.
Then the sequences of circles that they meet cannot be the same, so~neither can the semi-infinite paths in $T_\R$.

This argument also shows that, given an end point $e$ that corresponds to a sequence $(C_0, C_1, \dots)$ of circles, the set $\{ B_k \}_{k \in \mathbf{N}}$ (each $B_k$ as described above) is a basis of neighborhoods at $e$.
\end{remark}

\subsection{Legal coloring of the tree of circles}

In this and the following sections, we~focus our attention on the regular rabbits, which include the basilica and the Douady rabbit.

\begin{definition}
\label{def.biregular.tree}
Any tree $T$ has a unique natural bipartition $\{A,B\}$ of vertices in which any two vertices lie in the same part of the partition if and only if their distance is even. A~tree $T$ is $(n,m)$\emph{-biregular} if every vertex of $A$ has degree $n$ and every vertex of $B$ has degree $m$, where $n, m \in \mathbf{N} \cup \{\infty\}$.
\end{definition}

\begin{remark}
For every rabbit $\R$, its tree of circles $T_\R$ is bipartitioned into the sets of circles and of cut points of $\R$. By~\cref{lem.tree.of.circles.is.tree}, $T_\R$ is an $(n,\infty)$-biregular tree if and only if $\R$ is an $n$-regular rabbit.
\end{remark}

For each circle $C$ of a rabbit, we~want to encode the cyclic structure of the cut points belonging to $C$ in the tree $T_\R$. We~will do this by using certain colorings of~$T_\R$. We~cannot employ the renowned construction of \cite{BM00} because the trees of circles are not regular trees of finite degree. We~instead use the technology developed in \cite{biregular}, which generalizes the construction of Burger and Mozes to the case of biregular trees and allows infinite degree.
Let us first describe how this works.

\begin{definition}
\label{def.half.edges}
Given a graph $T$ (undirected and without loops nor parallel edges), a \emph{half-edge} of $T$ is an ordered pair of adjacent vertices. If~$T$ is a graph, we~denote by $H(T)$ the set of half-edges of $T$.
Given a vertex $v$ of a graph, we~denote by $\mathrm{out}(v)$ the set of half-edges that originate from $v$ and by $\mathrm{in}(v)$ the set of half-edges that terminate at $v$.
\end{definition}

In \cite{biregular} half-edges are instead called \textit{arcs}, but we avoid this term because we are already using it for the arcs of topological spaces.
The term half-edge is inspired by the fact that each edge between $v$ and $w$ corresponds precisely to two half-edges, namely $(v,w)$ and $(w,v)$.

\begin{remark}
\label{rmk.rabbit.tree.separation.relation}
Let $T_\R$ be the tree of circles of a rabbit. If~$v = C$ is a circle, then we can naturally identify $\mathrm{out}(C)$ with the set of cut points that belong to $C$. If~instead $v = p$ is a cut point, then we identify $\mathrm{out}(p)$ with the set of circles that $p$ belongs to, or equivalently with the set of connected components of $\R \setminus \{p\}$.

Each circle $C$ has two natural cyclic orders, each inducing the same separation relation (see \cref{sec.order.and.separation}).
Thanks to condition \cref{def.rabbit1} of \cref{def.rabbit}, restricting these relations to the set of cut points belonging to $C$ yields a countable dense cyclic order.
With the identification of $\mathrm{out}(C)$ with the set of cut points that belong to $C$, we~thus have two natural dense cyclic orders inducing a unique separation relation on $\mathrm{out}(C)$, for each circle $C$ of the rabbit $\R$
\end{remark}

In order to equip $T_\R$ with the dense cyclic order on $\mathrm{out}(C)$, we~rely on the tool defined right below.

\begin{definition}[\cite{biregular}]
\label{def.legal.coloring}
Let $T$ be an $(n,m)$-biregular tree ($n, m \in \mathbf{N} \cup \{\infty\}$) with bipartition $\{A,B\}$. A~\emph{legal coloring} of $T$ with colors $K_n$ and $K_m$ is a coloring $k$ of its half-edges, \ie a map
\[
 k \colon H(T) \to K_n \cup K_m, 
\]
that enjoys the following properties:
\begin{enumerate}
\item for all $v \in A$, $k$ maps $\mathrm{out}(v)$ bijectively to $K_n$;
\item for all $v \in B$, $k$ maps $\mathrm{out}(v)$ bijectively to $K_m$;
\item for all $v \in V(T)$, $k$ is constant on $\mathrm{in}(v)$.
\end{enumerate}
\end{definition}


Note that, by \cite[Prop.\,11]{biregular}, legal colorings form an orbit under the action of $\Aut(T)$, \ie every two legal colorings with the same set of colors differ from an automorphism of $T$ and the image of a legal coloring is another legal coloring.
The third condition of \cref{def.legal.coloring} is key to this.

\begin{lemma}
\label{lem.coloring.of.tree.of.circles}
Let $\R$ be an $n$-regular rabbit.
Let $\Omega$ be a countable set together with a dense cyclic order $O_\Omega$ and, if $n$ is finite, let $[n] = \{1, \dots, n\}$ be equipped with its standard cyclic order.
For each circle $C$ of $\R$, equip $\mathrm{out}(C)$ with a countable dense cyclic order on the cut points of $C$ as described in \cref{rmk.rabbit.tree.separation.relation} and, if $n$ is finite, for each cut point $p$ of $\R$, equip $\mathrm{out}(p)$ with a cyclic order.
Then there is a legal coloring with colors $([n],\Omega)$ whose restriction to $\mathrm{out}(v)$ is an isomorphism of cyclic orders for each vertex $v$ (only for those that correspond to circles, if $n$ is infinite). In~particular, by \cref{lem.preserves.respects.cyclic.order} it is an isomorphism of the induced separation relations (\cref{def.separation.relation}).
\end{lemma}

\begin{proof}
We are going to repeatedly use the fact that the set of cut points of each circle $C$ is naturally endowed with two countable dense cyclic orders that induce the same separation relation.
On each circle, we~fix any of the two.
This induces a unique separation relation on each $\mathrm{out}(C)$, as described in \cref{rmk.rabbit.tree.separation.relation}.

Fix a starting circle $C_0$ of $\R$.
The correspondence between the set $\mathrm{out}(C_0)$ of half-edges of $T_\R$ originating at $C_0$ and the set of cut points of $C_0$ induces the desired bijection $\mathrm{out}(C) \to \Omega$ that preserves the cyclic orders.

Proceeding by induction, suppose that $T$ is a subtree of $T_\R$ whose leaves are circles of $\R$ and that we have already built the map $k$ on the half-edges originating from $T$ (\ie on the half-edges $(v,w) \in \mathrm{out}(v)$ for each $v \in V(T)$, even if $w$ does not belong to $T$) in such a way that
\begin{enumerate}
\item\label{enumb:1} for each circle $C$ of $T$, $k$ maps $\mathrm{out}(C)$ bijectively to $\Omega$ and preserves the cyclic orders;
\item\label{enumb:2} for each cut point $p$ of $T$, $k$ maps $\mathrm{out}(p)$ bijectively to $[n]$ and, if $n$ is finite, preserves the cyclic orders;
\item\label{enumb:3} for each vertex $v$ of $T$, $k$ is constant on $\mathrm{in}(v)$.
\end{enumerate}
Suppose that $p$ is a cut point that does not belong to $T$ and is adjacent in $T_\R$ to some circle $C$ of $T$. We~need to show that we can add $p$ to $T$ along with all circles that are adjacent to it and extend the coloring in such a way that the above conditions are satisfied.
This will eventually color every vertex of $T_\R$.
Given $p$ and $C$ as above, define $k(p,C)$ as $k(q,C)$ for any other cut point $q$ in $T$ that belongs to $C$ (in the first step, when $V(T) = \{C\}$, this choice is instead arbitrary).
Then complete $k(p,-)$ to a bijection such that $k(p,C')$ satisfies condition \cref{enumb:2}.
Next, for each circle $C' \neq C$ that is adjacent to $p$ define $k(C',p) = k(C,p)$, so~that condition \cref{enumb:3} holds at the vertex $p$.
For each $C'$, color $\mathrm{out}(C')$ in such a way that condition \cref{enumb:1} holds.
The map $k$ extended in this way satisfies all conditions, so~we are done.
\end{proof}

In the remainder of this section, which of the two cyclic orders we fix on each circle will not matter, as they induce the same separation relation.
The cyclic order on $\mathrm{out}(p)$ (for $p$ a cut point) too is not going to matter here.


\subsection{Universal groups of biregular trees}
\label{sub.universal.groups}

A graph automorphism $g$ permutes the half-edges around each vertex. In~particular, for any vertex $v$ we consider the bijection $g|_{\mathrm{out}(v)} \colon \mathrm{out}(v) \to \mathrm{out}(g(v))$.

\begin{definition}
\label{def.universal.group}
Let $T$ be an $(n,m)$-biregular tree with bipartition $\{A, B\}$ and fix a legal coloring $k$ of $T$.
Given $N \leq \Sym(K_n)$ and $M \leq \Sym(K_m)$, the associated \emph{universal group} with prescribed local actions $\U_k (N,M)$ (or simply $\U (N,M)$ if $k$ is understood) is the group of those automorphisms $g$ of $T$ that preserve the bipartition and such that, for all $v \in V(T)$,\vspace*{-3pt}
\begin{equation}
k|_{\mathrm{out}(g(v))} \circ g|_{\mathrm{out}(v)} \circ k|_{\mathrm{out}(v)}^{-1} \in
\begin{cases}
N & \text{if $v \in A$},\\
M & \text{if $v \in B$},
\end{cases}
\end{equation}
\ie the permutation of colors induced by $g$ at the vertex $v$ belongs to $N$ or $M$, depending on whether $v$ belongs to $A$ or $B$. We~endow $\U_k(N,M)$ with the pointwise convergence topology.
\end{definition}

We will denote the color permutation induced by $g$ around $v$ by\vspace*{-3pt}
\[
 \sigma_g(v) \coloneqq k|_{\mathrm{out}(g(v))} \circ g|_{\mathrm{out}(v)} \circ k|_{\mathrm{out}(v)}^{-1}. 
\]
This is called the \emph{local action} of $g$ at $v$.

The paper \cite{biregular} studies the topological and transitive properties that $\U(N,M)$ inherits from $N$ and $M$. It~also shows that $\U(N,M)$ is topologically isomorphic to its permutation group of the vertices of one of the two parts of the bipartition, which Smith calls the \textit{box product} and denotes by $N \boxtimes M$.

\begin{remark}
\label{rmk.BM.and.biregular}
As noted in the third paragraph of Remark 4 of \cite{biregular}, when $n=2$ there is natural identification between the universal group $\U(\Sym([2]),M)$ and the Burger-Mozes group $\U(M)$ that acts on an $m$-regular tree.
More precisely, $\U(\Sym([2]),M)$ and $\U(M)$ are permutation isomorphic if we consider the action of the first on the set of vertices of degree $m$.
This identification is found simply by taking the barycentric subdivision of the $m$-regular tree.

In particular, it~is not hard to see that the group $\U(\Sym([2]),\Aut(S))$ that we will describe in the next subsection for the basilica (which is the $2$-regular rabbit) is the same as the Burger-Mozes group found in \cite{Neretin}.
\end{remark}

\subsection{Uniqueness of the regular rabbits}


The following proposition shows that there is a unique $n$-regular rabbit for each $n$ (up to homeomorphisms).
Given a rabbit $\R$ and its tree of circles $T$, below we denote by $\pi \colon \R \to \overline{T}$ the map that sends each end point to the corresponding point of $\partial T$, each cut point to itself (as a vertex of $T$) and each regular point to the unique circle to which it belongs (as a vertex of $T$).

\begin{proposition}
\label{prop.rabbit.surjectivity}
Let $\R_1$ and $\R_2$ be two $n$-regular rabbits and let $T$ be their tree of circles (which is the same up to isomorphisms by \cref{lem.tree.of.circles.is.tree}).
Let $k_1$ and $k_2$ be two legal colorings of $T$ as in \cref{lem.coloring.of.tree.of.circles}.
Then, for every automorphism $\varphi$ of $T$ whose local action on $\mathrm{out}(v)$ at each vertex $v$ is an isomorphism between the two separation relations, there is a homeomorphism $\psi \colon \R_1 \to \R_2$ that makes the following diagram commute.\vspace*{-3pt}
\begin{center}
\begin{tikzcd}
\R_1 \arrow[r, "\ts\psi"] \arrow[d, "\ts\pi_1"] & \R_2 \arrow[d, "\ts\pi_2" ] \\
T \arrow[r, "\ts\varphi"] & T
\end{tikzcd}
\end{center}
\end{proposition}

\begin{proof}
Let $\R_1$ and $\R_2$ be two $n$-regular rabbits and let $T_1$ and $T_2$ be their trees of circles.
Fix two legal colorings satisfying \cref{lem.coloring.of.tree.of.circles} and let $\varphi \colon T_1 \to T_2$ be a color-preserving automorphism (the local action at any vertex is the identity).
Note that the maps $\pi_1$ and $\pi_2$ are injective on the sets of cut points of $\R_1$ and $\R_2$.
This defines a bijection $\psi = \pi^{-1}_2 \circ \varphi \circ \pi_1$ between the cut points of $\R_1$ and $\R_2$.
Since $\varphi$ is a tree isomorphism and the edge adjacency in $T_1$ and $T_2$ is given by the inclusion of cut points to circles, the map $\psi$ restricts to bijections between the cut points of each circle.
Moreover, since $\varphi$ is color-preserving, the circular order of cut points on each circle is preserved, and since the sets of cut points are dense in $\R_1$ and $\R_2$ (by~condition~\cref{def.rabbit1} of \cref{def.rabbit} together with \cref{lem.coloring.of.tree.of.circles}), $\psi$ extends uniquely to a continuous map between a circle of $\R_1$ and its image.
Doing this for all circles of $\R_1$ extends $\psi$ to a bijection between $\R_1$ and $\R_2$ such that $\varphi \circ \pi_1 = \pi_2 \circ \psi$. It~now remains to show that~$\psi$ is a homeomorphism.

Note that the connected components of $\R_i$ minus finitely many circles correspond precisely to the connected components of $T_i$ minus finitely many circles together with their cut points (seen as vertices of the tree).
Then, if $U$ is such a subset of $\R_2$, its preimage $\psi^{-1}(U)$ is a connected component of $\R_1$ minus finitely many circles, so~in particular it is open. By~\cref{cor.rabbit.subbasis}, the collection of such subsets is a subbasis for the topology of a rabbit, so~this shows that $\psi$ is continuous.
Ultimately, this shows that $\psi$ is a continuous bijection from a compact space to a Hausdorff space, so~it is a homeomorphism.
\end{proof}

By \cref{lem.coloring.of.tree.of.circles}, we~can always find legal colorings of any two $n$-regular rabbits satisfying the hypotheses of \cref{prop.rabbit.surjectivity}.
Hence we obtain the desired result.

\begin{theorem}
\label{thm.uniqueness.rabbit}
Any two $n$-regular rabbits are homeomorphic.
\end{theorem}

Thanks to this theorem, we~can now refer to any $n$-regular rabbit as \textit{the} $n$-regular rabbit.
For example, \textit{the} basilica and \textit{the} Douady rabbit.


\subsection{Homeomorphism groups of regular rabbits}
\label{sub.rabbit.homeo.group.is.universal}

Here we show that the group of homeomorphisms of a regular rabbit is the universal group $\U(\Sym([n]),\Aut(S))$, where $S$ is the separation relation of the countable dense cyclic order.

\begin{proposition}
\label{prop.homeo.rabbit.acts.on.tree}
Given an $n$-regular rabbit $\R$, the group $\Homeo(\R)$ acts faithfully by graph automorphisms on $T_\R$ in such a way that defines an embedding
\[
 \Pi \colon \Homeo(\R) \hto \U_k (\Sym([n]),\Aut(S)), 
\]
where $k$ is the legal coloring of $T_\R$ provided by \cref{lem.coloring.of.tree.of.circles}
\end{proposition}

\begin{proof}
A homeomorphism $\varphi$ of $\R$ induces a permutation of the circles and of the cut points, which are the vertices of the tree of circles $T_\R$.
Recall that $\{p,C\}$ is an edge of $T_\R$ if and only if $p \in C$, so~clearly the permutation induced by $\varphi$ preserves the edge adjacency of $T_\R$.
Thus, $\Homeo(\R)$ acts by graph automorphisms on $T_\R$.
The set of cut points is dense in $\R$ by condition \cref{def.rabbit1} of \cref{def.rabbit} and \cref{cor.rabbit.cut.points}, so~the action is faithful.

Now, consider the legal coloring $k$ of $T_\R$ provided by \cref{lem.coloring.of.tree.of.circles}.
The bipartition of $T_\R$ distinguishes between cut points and circles, so~clearly the action of $\Homeo(\R)$ preserves it.
For each circle $C$ of $\R$, every homeomorphism $\varphi$ of $\R$ maps bijectively between the sets of cut points belonging to $C$ and those belonging to $\varphi(C)$ in a separation-preserving fashion.
This means that the permutation of colors induced by~$\varphi$ around the vertex $C$ of $T_\R$ belongs to $\Aut(S)$, as needed.
\end{proof}

The surjectivity of this embedding follows immediately from \cref{prop.rabbit.surjectivity}, so~we have the following fact.

\begin{proposition}
\label{prop.rabbit.embedding.is.surjective}
The embedding $\Pi$ of \cref{prop.homeo.rabbit.acts.on.tree} is surjective.
\end{proposition}

We are now only missing the bicontinuity of $\Pi$. We~endow $\Homeo(\R)$ with the uniform convergence topology.

\begin{proposition}
\label{prop.rabbit.embedding.is.continuous}
The embedding $\Pi$ of \cref{prop.homeo.rabbit.acts.on.tree} is a homeomorphism.
\end{proposition}

\begin{proof}
Since both are Polish groups, it~suffices to prove that the abstract group isomorphism $\Pi \colon \Homeo(\R) \to \U(\Sym([n]),\Aut(S))$ is continuous (for example because of \cite[Cor.\,2.3.4]{MR2455198}).
Suppose that $g_m$ is a sequence in $\Homeo(\R)$ that converges to the identity and consider an arbitrary circle $C_0$ of the rabbit $\R$. We~will show that, for large enough $m \in \mathbf{N}$, the homeomorphism $g_m$ fixes $C_0$.
Let $C_1$, $C_2$ and~$C_3$ be circles of $\R$ that lie in three distinct connected components $U_1$, $U_2$ and $U_3$ of $\R \setminus C_0$.
Then, for $m$ large enough, say $m \geq N$, the circle $g_m (C_i)$ is included in $U_i$ for each $i = 1, 2, 3$. In~the tree of circles, consider the unique paths between $C_1$ and $C_2$, between~$C_1$ and~$C_3$ and between $C_2$ and $C_3$:
by \cref{rmk.arcs.and.paths}, they intersect precisely at $C_0$, so~the vertex $C_0$ is fixed by $\Pi(g_m)$ for all $m \geq N$.
This argument shows that each vertex is fixed by $\Pi(g_m)$ for large enough $m$, so~the sequence $\Pi(g_m)$ converges to the identity of $\U(\Sym([n]),\Aut(S))$ and thus $\Pi$ is continuous.
\end{proof}

Together, the previous propositions immediately imply the desired result.


\begin{theorem}
\label{thm.homeo.rabbits.are.auto.trees}
The group $\Homeo(\R_n)$ of homeomorphisms of an $n$-regular rabbit~$\R_n$ is topologically isomorphic to $\U (\Sym([n]),\Aut(S))$.
\end{theorem}

\setcounter{theorem}{30}
\begin{corollary}
\label{cor.rabbit.groups.closed}
For every $n \geq 2$, the group $\Homeo(\R_n)$ embeds into $\Aut(T_{n,\infty})$ as a closed subgroup.
\end{corollary}

\begin{proof}
\cite[Th.\,6(v)]{biregular} states that $\U(N,M)$ is closed in $\Aut(T)$ if~$N$ and~$M$ are closed in $\Sym([n])$ and $\Sym([m])$, respectively.
Since $\Aut(S)$ is a closed subgroup of the infinite symmetric group (\cref{rmk.Aut.S.C.closed}), we~are done.
\end{proof}

\appendix
\section{The countable dense cyclic order and its separation relation}
\label{sec.order.and.separation}

In this section we gather some useful facts about the dense cyclic order $O$ and its separation relation $S$ that we use throughout the paper.

\Subsection{Cyclic orders and their separation relations}

\begin{definition}[Cyclic order]
\label{def.cyclic.order}
Let $\Omega$ be a set. A~\emph{cyclic order} on $\Omega$ is a ternary relation $[\ ,\ ,\ ]$ on $X$ that enjoys the following properties:
\begin{enumerate}
\item \emph{Cyclicity:} If $[a, b, c]$ then $[b, c, a]$,
\item\emph{Asymmetry:} If $[a, b, c]$ then not $[c, b, a]$,
\item\emph{Transitivity:} If $[a, b, c]$ and $[a, c, d]$ then $[a, b, d]$
\item\emph{Connectedness:} If $a, b,c$ are distinct, then either $[a, b, c]$ or $[c, b, a]$.
\end{enumerate}
\end{definition}

Given a cyclic order, one can define the inverse cyclic order $[x,y,z]^{-1}\iff[z,y,x]$. It~is readily seen that this is a well-defined cyclic order.

Let us denote by $\Omega^{(3)}$ the set of all triples $(x,y,z) \in \Omega^3$ such that $x,y,z$ are distinct.
Every cyclic order defines an \emph{orientation map}
\[
 O \colon \Omega^{(3)} \to \{+,-\}, 
\]
that is $O(x,y,z)=+ \iff [x,y,z]$.
Since a cyclic order is entirely determined by its orientation map, we~take the liberty of using the symbol $O$ to also refer to the cyclic order.

Following Vailati \cite[\S204, p.\,215]{Russel}, we~introduce the following quaternary relation (see Coxeter \cite[Ch.\,3]{Coxeter}, for a more recent reference).

\begin{definition}[Separation relation]
\label{def.separation.relation}
Let $\Omega$ be a set equipped with a cyclic order $[\ ,\ ,\ ]$.
The \emph{separation relation} associated to it is the relation $S$ such that $S(a,b,c,d) \iff [a,b,c]\wedge [c,d,a]$ or $[a,d,c]\wedge [c,b,a]$.
\end{definition}

It is easy to verify that a cyclic order and its inverse induce the same separation relation.

\begin{example}
Let $\Omega=\mathbf{S}^1$ be the unit circle in $\mathbf{C}$ with the following cyclic order.
For $z_0,z_1,z_2\in \mathbf{S}^1$, choose two preimages $x_1,x_2\in\mathbf{R}$ of $z_1,z_2$ via the exponential map $x\mto \exp(ix)$ that lies in the same connected component of $\mathbf{R}\setminus\exp^{-1}(z_0)$ and write $[z_0,z_1,z_2]\iff x_1<x_2$ for the usual order on $\mathbf{R}$.
\end{example}

\begin{example}
The separation relation associated to the cyclic order on $\mathbf{S}^1$ can be described in a topological way:
$S(a,b,c,d)$ holds if and only if $b$ and $d$ lie in different connected components of $\mathbf{S}^1\setminus\{a,c\}$.
That is, $\{b,d\}$ is separated by $\{a,c\}$ and vice versa.
\end{example}

Linear orders, cyclic orders, betweenness and separations relations are closely related, as explained in \cite{Huntington}. In~particular, given a cyclic order on a set $\Omega$ and $x\in\Omega$, we~can set $y<_xz$ if $(x,y,z)$ is cyclically ordered.
The binary relation $<_x$ is a linear order on $\Omega$.
Thus, given a subset $F$ of a cyclically ordered set $\Omega$, we~call \emph{cyclic successor} of $x\in F$ the minimum of the linear order $<_x$ on $F\setminus\{x\}$ (if it exists).
Conversely if $<$ is a linear order on some set $\Omega$, we~can define a cyclic order via $[x,y,z]\iff x<y<z\vee y<z<x \vee z<x<y$. In~particular, we~call \emph{standard} the cyclic order on $[n]=\{1,\dots,n\}$ induced by the natural linear order $<$.

Let $\Omega$ be a set with a cyclic order $[\ ,\ ,\ ]$. We~say that a bijection $g$ of $\Omega$ \emph{preserves} the cyclic order if $[x,y,z]\iff [g(x),g(y),g(z)]$ and that it \emph{reverses} the order if $[x,y,z]\iff [g(x),g(y),g(z)]^{-1}$. If~$S$ is a separation relation, we~say that a bijection~$g$ of~$\Omega$ \emph{preserves} $S$ if $S(x,y,z,w)\iff S(g(x),g(y),g(z),g(w))$.

\begin{lemma}
\label{lem.preserves.respects.cyclic.order}
Let $\Omega$ be a set with a cyclic order and let $S$ be the associated separation relation. A~bijection $g$ of $\Omega$ preserves $S$ if and only if it preserves or reverses the cyclic order.
\end{lemma}

\begin{proof}
Since the cyclic order and its inverse cyclic order induce the same separation relation, it~is clear that a bijection that preserves or reserves the cyclic order preserves the separation relation.

Conversely, assume that $(u,v,w)$ satisfies $[u,v,w]$.
Then, for any triple $a,b,c$,
\begin{align*}
[a,b,c]&\iff (S(a,u,v,w)\wedge S(a,b,v,w)\wedge S(a,b,c,w))\\
&\hspace*{2cm}\vee(S(b,u,v,w)\wedge S(b,c,v,w)\wedge S(b,c,a,w))\\
&\hspace*{2cm}\vee(S(c,u,v,w)\wedge S(c,a,v,w)\wedge S(c,a,b,w))\\
&\hspace*{2cm}\vee(S(a,u,v,w)\wedge S(a,b,v,w)\wedge S(a,b,w,c))\\
&\hspace*{2cm}\vee(S(b,u,v,w)\wedge S(b,c,v,w)\wedge S(b,c,w,a))\\
&\hspace*{2cm}\vee(S(c,u,v,w)\wedge S(c,a,v,w)\wedge S(c,a,w,b))
\end{align*}
Let us write $R(a,b,c,u,v,w)$ for the relation with $6$ variables that appears on the right-hand side of the equivalence above.
Let us fix $u,v,w \in \Omega$ such that $[u,v,w]$ and let us assume that $g$ preserves $S$. If~$[g(u),g(v),g(w)]$ then $[a,b,c]$ implies $R(a,b,c,u,v,w)$ and, since $g$ preserves $S$, $R(g(a),g(b),g(c),g(u),g(v),g(w))$ and thus $[g(a),g(b),g(c)]$.
Similarly if $[g(u),g(v),g(w)]^{-1}$, for $a,b,c$ such that $[a,b,c]$, $R(a,b,c,u,v,w)$ holds and thus $[g(a),g(b),g(c)]^{-1}$.
This means that $g$ preserves or reverses the cyclic order.
\end{proof}

\subsection{Countable dense cyclic orders}

A cyclic order is \emph{dense} if for all distinct $x,z\in \Omega$, there is $y \in \Omega$ such that $[x,y,z]$.
There is a unique countable dense cyclic order $O$ up to isomorphism (this follows from the analogous statement for linear orders).
Such a cyclic order $O$ has been studied in \cite{Truss}.
The Dedekind completion of $O$ is $\mathbf{S}^1$ with its usual topology.

\begin{example}
A realization of this unique dense countable cyclic order can be obtained by choosing any dense countable subset of the circle $\mathbf{S}^1$.
For example, if we identify $\mathbf{S}^1$ with $\mathbf{R}/\mathbf{Z}$ then $\mathbf{Q}/\mathbf{Z}$ and $\mathbf{Z}[1/2]/\mathbf{Z}$ are standard choices.
\end{example}

Throughout the rest of this appendix and all other sections, we~denote by $O$ this countable dense cyclic order and by $S$ the associated separation relation. We~denote by $\Aut(S)$ the group of all bijections of $\Omega$ that preserve this separation relation $S$ and by $\Aut(O)$ the subgroup of $\Aut(S)$ of bijections of $\Omega$ that preserve this cyclic order~$O$.
Let us gather a few facts about these two groups.

\begin{remark}
\label{rmk.Aut.S.C.closed}
The structures $(\Omega,O)$ and $(\Omega,S)$ are Fraïssé limits of all finite cyclic orders and all finite separation relations respectively. In~particular, it~follows that $\Aut(S)$ and $\Aut(O)$ are closed subgroups of $\Sym([\infty])$ for the pointwise convergence topology.
\end{remark}

\begin{thebibliography}{99}
\bibitem[Nad92]{continua} S.B.Nadler,Jr., \emph{Continuum theory}, Monographs and Textbooks in Pure and Applied Math.158, Marcel Dekker,Inc.,New York,1992.
\bibitem[BM00]{BM00} M.Burger and S.Mozes, ``Groups acting on trees:from local to global structure'', \emph{Publ. Math. Inst. Hautes \'Etudes Sci.}\textbf{92}(2000),113--150.
\bibitem[Smi17]{biregular} S.M.Smith, ``A product for permutation groups and topological groups'', \emph{Duke Math. J.}\textbf{166}(2017),no.15,2965--2999.
\bibitem[Ner24]{Neretin} Y.A.Neretin, ``On the group of homeomorphisms of the basilica'', \emph{Topology Proc.}\textbf{64}(2024),121--128.
\bibitem[Gao09]{MR2455198} S.Gao, \emph{Invariant descriptive set theory}, Pure and Applied Math.293, CRC Press,Boca Raton,FL,2009.
\bibitem[Rus03]{Russel} B.Russell, \emph{The principles of mathematics.Vol.I}, Cambridge University Press,Cambridge,1903.
\bibitem[Cox93]{Coxeter} H.S.M.Coxeter, \emph{The real projective plane}, third ed., Springer-Verlag,New York,1993.
\bibitem[Hun35]{Huntington} E.V.Huntington, ``Inter-relations among the four principal types of order'', \emph{Trans. Amer. Math. Soc.}\textbf{38}(1935),no.1,1--9.
\bibitem[Tru09]{Truss} J.K.Truss, ``On the automorphism group of the countable dense circular order'', \emph{Fund. Math.}\textbf{204}(2009),no.2,97--111.
\end{thebibliography}
\end{document}
